\section{Effective support bounds and exact feasibility}\label{app:effective}
This appendix records the analytic estimates used by the finite interval
calculations. All constants below are evaluated on whole parameter and root
boxes. Point evaluations at approximate roots do not suffice.

\subsection{Local estimates}
Put \(a_0=2^{-14}\), \(\epsilon=2^{-12}\) and
\(b(a)=b^*+va^2\). At an original residual zero \(z_k\), start with the
center \(c_k=z_k+\kappa_ka^2\), with \(\kappa_k\) from \eqref{eq:kappa}.
Write \(K=|\kappa_k|\),
\(Q_{j,\ell,k}=\sup_{N_k}|q_j^{(\ell)}|\), and
\(R_{\ell,k}=\sup_{N_k}|q_*^{(\ell)}|\).
In the following formulas, unlabelled derivatives are evaluated at \(z_k\).
The moment change at this switch is exactly
\[
 \Delta_{j,k}=-2\sigma_k\mathbb E
                 \int_0^{\kappa_ka^2+ay}q_j(z_k+x)\dd x .
\]
Taylor's formula gives
\[
 |\Delta_{j,k}-a^2(-2\sigma_kq_j\kappa_k-\sigma_kq_j'/3)|
                  \leq H_{j,k}a^4 ,
\]
where
\begin{align*}
 H_{j,k}={}&|q_j'|K^2
 +\frac{|q_j''|}{3}(K+K^3a_0^2)\\
 &+\frac{Q_{j,3,k}}{12}(1/5+2K^2a_0^2+K^4a_0^4).
\end{align*}
For the residual, whose value at its zero is exactly zero,
\[
 \Delta_{*,k}=-\sigma_kr_ka^2/3
 -\sigma_k(r_k\kappa_k^2+t_k\kappa_k/3+u_k/60)a^4+\varepsilon_k,
 \qquad |\varepsilon_k|\leq K_{*,k}a^6,
\]
with
\begin{align*}
 K_{*,k}={}&|t_k|K^3/3+|u_k|(2K^2+K^4a_0^2)/12\\
 &+|q_*^{(4)}|(K+(10/3)K^3a_0^2+K^5a_0^4)/60\\
 &+R_{5,k}(1/7+3K^2a_0^2+5K^4a_0^4+K^6a_0^6)/360 .
\end{align*}
These coefficients come from signed averages of powers, including
\(\mathbb E(\kappa a^2+ay)^5=\kappa a^6+(10/3)\kappa^3a^8+\kappa^5a^{10}\).
Bounding that odd average by an absolute fifth moment would lose the
required order.

For \(F_k(c)=\mathbb E q_{b(a)}(c+ay)\), the quadratic coefficient at
\(c_k\) vanishes because \(r_k\kappa_k+t_k/6-Q_k^{\mathsf T}v=0\).
Thus \(|F_k(c_k)|\leq h_ka^4\), where
\begin{align*}
 h_k={}&|t_k|K^2/2+|u_k|(K+K^3a_0^2)/6\\
 &+R_{4,k}(1/5+2K^2a_0^2+K^4a_0^4)/24\\
 &+\sum_{j<3}|v_j|\{
       |q_j'|K+Q_{j,2,k}(1/3+K^2a_0^2)/2\}.
\end{align*}
If the signed derivative \(\sigma_kF_k'\geq m_k>0\) on the relevant
center interval, the true balanced center is within \(h_ka^4/m_k\).
The local objective has second derivative at most \(-2m_k\).
Its possible increase after balancing is consequently at most
\(h_k^2a^8/m_k\). Containment of both ramps is checked before this
estimate is applied.

Move the first three centers by \(a^4y_k\), with
\(|y_k|\leq W_k\), \(W=(2^{26},2^{26},2^{24})\).
For \(J_{jk}=-2\sigma_kq_j(z_k)\), let \(C\) be the dyadic
preconditioner recorded for the cell. The Jacobian deviation is bounded by
\[
 \Delta J_{jk}\leq
 2\{Q_{j,1,k}(K_ka_0^2+W_ka_0^4)+Q_{j,2,k}a_0^2/6\}.
\]
In the norm \(\max|y_k|/W_k\), the map subtracting \(C\) times the
lower moments divided by \(a^4\) has the following uniform bounds:
\[
 \text{derivative rows}<(.005731,.007842,.004703),\qquad
 \text{forcing}<(.328781,.448496,.267009).
\]
Each corresponding sum is less than one. The contraction theorem gives
exactly vanishing lower moments within this center construction.
The same estimates imply that \(C\) is nonsingular.
With \(\mathcal L_k=\sup|q_{b(a)}'|\), the repair costs at most
\[
 a^8\sum_{k=1}^3(2h_kW_k+\mathcal L_kW_k^2).
\]
This estimate is made for \(q_{b(a)}\), whose center derivative is of
order \(a^4\). The repaired target equals its \(q_3\) target because
the lower moments vanish.

\subsection{A changing finite prefix and its entire omitted tail}
Always retain the 28 zeros below \(1\). Retain a tail zero \(z_k>1\)
precisely when \(ae^{4z_k}\leq\epsilon\).
This gives a finite prefix for every \(a>0\).
After the last active ramp, extend its final plateau to infinity.
The first inactive zero is separated from the finite interval by the
validated positive gap \([1,1.001]\).

Uniformly on the driver cells, \(|v|<(5,4,180)\) and the expanded
trial multiplier box satisfies
\(|b|<(.001,.1,.005)\).
On each active tail neighborhood \([z_k-3a,z_k+3a]\), the raw phase
estimates give \(|\kappa_k|<46e^{4z_k}\), signed raw trial derivative
greater than \(550z_k^3\), and raw absolute value less than \(2800az_k^3\).
Also \(\tfrac12<w(x)/w(z_k)<2\), since the logarithmic variation
is less than \(408\epsilon<1/2\).
Combining the weight derivative with these estimates gives
\(\sigma_k q_{b(a)}'>200w(z_k)z_k^3\).
Thus balancing is valid on active tail neighborhoods as well.

For \(u\geq1\), the weight satisfies
\[
 w_\nu(u)\leq88e^{9u+9u^4-\pi e^{4u}} .
\]
Writing \(P_\nu=\lambda u^2+\mu u^4+\nu u^6\) and \(H=1/64\),
\begin{align*}
 |P_\nu'|&\leq44u^3+6Hu^5,&
 |P_\nu''|&\leq116u^2+30Hu^4,\\
 |P_\nu'''|&\leq216u+120Hu^3,&
 |P_\nu^{(4)}|&\leq216+360Hu^2,&
 |P_\nu^{(5)}|&\leq720Hu .
\end{align*}
The inequalities \(e^{4u}>50u^3\) and \(u^5e^{-4u}<1/40\) imply
\(|P_\nu^{(j)}|<e^{4ju}\), \(1\leq j\leq5\).
The theta derivatives follow from
\[
 P_0(x)=-3+2x,\qquad
 P_{\ell+1}(x)=(5-4x)P_\ell(x)+4xP_\ell'(x),
\]
applied term by term to \(\Phi\); product rules and exponential
derivative polynomials give the needed five weight derivatives.
This retains the sextic terms in local calculations and discards
their negative exponent only when taking an upper tail bound.

For \(0\leq p\leq3\), \(d\in\{12,16,24\}\), the corresponding root
sum and integral are bounded by
\[
 S_d=\frac{88e^{18+d-\pi e^4}}{1-50^{-4}},\qquad
 I_d=\frac{88}{500}e^{18+d-\pi e^4}.
\]
Indeed the logarithmic derivative of the envelope is below \(-500\);
root spacing exceeds \(3/85\), giving a successive term ratio below
\(50^{-4}\).
For every nonnegative sequence \(t_k\), inactive roots satisfy
\[
 \sum_{ae^{4z_k}>\epsilon}t_k
       \leq(a/\epsilon)^s\sum_k t_ke^{4sz_k}\quad(s>0).
\]
This supplies the missing powers of \(a\); a fixed cutoff would not.
For example the inactive quadratic and fourth coefficients are paid for
by \(300\epsilon^{-4}S_{16}\) and
\(2071000\epsilon^{-2}S_{16}\), respectively, at sixth order.
The lower-moment coefficient uses \(188\epsilon^{-2}S_{12}\).
Cutting the final integral at the first inactive zero minus \(3a\)
gives the additional bounds
\[
 \frac{18\epsilon^{-6}I_{24}}{1-72a_0},\qquad
 \frac{2^{j+1}\epsilon^{-4}I_{16}}{1-48a_0}
\]
for the normalized objective and lower-moment tails.
Active tail derivatives use the same envelope with the local formulas
above. The evidence records these sums, signs, and containment margins.

Combining the finite part, the active tail, and the inactive tail yields,
for every fixed \((\nu,a)\in I\times(0,a_0]\),
\begin{align}
 H_{b(a)}(a)&\leq D-\Gamma a^2/3+\Xi a^4+K_da^6,\label{eq:effectivesupport}\\
 |S(a)-(D-\Gamma a^2/3+\Xi a^4)|&\leq K_pa^6,\label{eq:effectiveprimal}
\end{align}
where \(K_d<655000\) and \(K_p<4.746\cdot10^6\) suffice uniformly.
During each three-center contraction, \(a\) and hence the prefix are
fixed. Continuity of the prefix as a function of \(a\) is not used.

\subsection{Scalar feasibility and inversion}
Set \(U=9.181\cdot10^{-10}\) and define
\[
 F_M(A)=DA-\Gamma A^3/(3M^2)+\Xi A^5/M^4,\qquad
 J_M(A)=F_M(A)-K_pA^7/M^6.
\]
On every driver cell, interval evaluation proves, for all \(M\geq M_0\),
\[
 J_M(\delta_0)<f<J_M(U),\qquad J_M'>0\text{ on }[\delta_0,U],
 \qquad U/M_0<a_0.
\]
The upper endpoint's normalized margin exceeds \(1.7683\cdot10^{-15}\).
Let \(A_+\) be the scalar root \(J_M(A_+)=f\).
At \(a=A_+/M\), \eqref{eq:effectiveprimal} gives \(S(a)\geq f/A_+\).
Multiplication of the repaired profile by \(A=f/S(a)\leq A_+\)
gives exact target and slope \(A/a\leq M\). This proves feasibility
at every budget despite changes in the prefix. Attainment follows as
in Theorem~\ref{thm:infinite}.

For the optimizer, \eqref{eq:effectivesupport} gives
\(F_M(\delta(M))\geq f-K_dU^7/M^6\).
The scalar upper root gives the corresponding upper comparison.
To bound the residual of \(P_4\), put
\[
 w=(\delta_0/M)^2,\quad t=\Gamma/(3D),\quad e=\Xi/D,\quad
 p=1+tw+(3t^2-e)w^2 .
\]
Then \(F_M(P_4)/f-1=p-twp^3+ew^2p^5-1\).
This is a polynomial of degree twelve whose coefficients through
degree two vanish exactly. Bounding all remaining coefficients on
\(0\leq w\leq(\delta_0/M_0)^2\) gives
\(|F_M(P_4)-f|\leq K_{\rm poly}/M^6\).
Also \(P_4\in[\delta_0,U]\) and
\[
 F_M'\geq d_{\min}=D_{\rm lo}-\Gamma_{\rm hi}U^2/M_0^2>0 .
\]
The inverse comparison constants are
\[
 K_-=\frac{K_{\rm poly}+K_dU^7}{d_{\min}}<1.001\cdot10^{-53},
 \qquad
 K_+=\frac{K_{\rm poly}+K_pU^7}{d_{\min}}<1.620\cdot10^{-53}.
\]
Rounding these bounds outwards gives \eqref{eq:thetaerror}.

\section{Verifying the transport condition on an unbounded domain}
\label{app:transport}
The uniform caps used here are
\[
 \lambda\in(-3.651,-3.645),\quad\mu\in(8.330,8.336),\quad
 \nu\in[-1/64,0],
\]
together with the derivative caps in the proof of
Corollary~\ref{cor:thetaorder}. For \(\Phi=\sum\phi_n\), put
\(x=\pi e^{4u}\) and, for \(n\geq2\),
\[
 R_n=\frac{\phi_n}{\phi_1}
 =n^2\frac{2n^2x-3}{2x-3}e^{-(n^2-1)x},\qquad
 \ell_n=9+\frac{12}{2n^2x-3}-4n^2x.
\]
Then \(\ell_n'=-16n^2x-96n^2x/(2n^2x-3)^2\).
For \(R=\sum_{n\geq2}R_n\),
\begin{equation}\label{eq:logtheta}
 L_\Phi=\ell_1+\frac{R'}{1+R},\qquad
 L_\Phi'=\ell_1'+\frac{R''}{1+R}
                      -\left(\frac{R'}{1+R}\right)^2 .
\end{equation}
On \(x\geq3\),
\[
 R_n\leq2n^4e^{-(n^2-1)x},\quad
 |R_n'|\leq10n^6xe^{-(n^2-1)x},\quad
 |R_n''|\leq100n^8x^2e^{-(n^2-1)x}.
\]
These follow by differentiating the displayed ratio; its logarithmic
derivative has absolute value at most \(5n^2x\), and the difference
of the logarithmic second derivatives has absolute value at most
\(20n^2x\). The summable majorants also justify the differentiations.

For \(0\leq u\leq1\), include \(n=2,\ldots,12\) explicitly.
The remaining errors in \(R,R',R''\) are bounded by
\[
 4\cdot13^4e^{-504},\quad60\cdot13^6e^{-504},\quad
 1800\cdot13^8e^{-504},
\]
because consecutive majorants have ratio below \(1/2\).
Arb evaluation of \eqref{eq:logtheta} on each of the 1024 whole
intervals \([i/1024,(i+1)/1024]\), using all parameter caps simultaneously,
then bounds
\begin{align*}
 L_\nu&=L_\Phi+2\lambda u+4\mu u^3+6\nu u^5,\\
 r&=-f'/f+4\kappa+\lambda'u^2+\mu'u^4+u^6+\kappa uL_\nu,\\
 r_u&=2\lambda'u+4\mu'u^3+6u^5+\kappa(L_\nu+uL_\nu').
\end{align*}
The largest recorded upper endpoint of \(r+\kappa+a_0|r_u|\) is
less than \(-0.02170697\); a separate outward rational reconstruction
proves the simpler bound \(<-0.0217<-1/50\) on every cell.

For \(u\geq1\), \(x>150\) and
\[
 R<64e^{-450},\quad |R'|<192000e^{-450},\quad
 |R''|<1152000000e^{-450}.
\]
Each is less than \(1/4\), and \(R'<0\). Equation~\eqref{eq:logtheta}
therefore gives
\[
 L_\Phi\leq10-4x,\quad |L_\Phi|\leq4x+1,\quad
 |L_\Phi'|\leq16x+2.
\]
For \(H=1/64\),
\begin{align*}
 L_\nu&\leq10-4x+36u^3<0,\\
 |L_\nu|&\leq4x+1+8u+36u^3+6Hu^5,\\
 |L_\nu'|&\leq16x+10+108u^2+30Hu^4 .
\end{align*}
Set \(k_-=.00259\), \(k_+=.00260\), \(l=.3264\),
\(m=.2942\), \(F=.0437\), and \(\bar k=k_--5a_0k_+>0\).
Using the sign of \(L_\nu\) in the term without absolute values,
\[
 r+\kappa+a_0|r_u|\leq P_6(u)-4\bar k\pi u e^{4u}=:Q(u),
\]
where, in increasing degree, the coefficients of \(P_6\) are
\[
 \begin{gathered}
 -F+5k_++a_0k_+,\quad 10k_-+a_0(2l+18k_+),\quad l,\\
 a_0(4m+144k_+),\quad m+36k_-,\quad
 a_0(6+36Hk_+),\quad1 .
 \end{gathered}
\]
All positive-degree coefficients are positive.
With \(S=\sum_{j=1}^6jP_j\), \(P_6'(u)\leq Su^5\).
The inequality \((1+4u)e^{4u}>270u^5\) for \(u\geq1\)
follows on writing \(z=u-1\), using \(e^4>109/2\) and the first six
Taylor terms for \(e^{4z}\). The difference polynomial has initial
coefficients \(5/2,-42,352\), forming a positive quadratic, and all
higher coefficients positive. Exact rational arithmetic then gives
\begin{align*}
 Q'(u)&<(S-4\bar k\cdot3\cdot270)u^5<0,\\
 Q(1)&<\sum P_j-4\bar k\cdot170\\
     &=-0.051116612396240234375<-1/50 .
\end{align*}
This proves the required inequality on the entire tail.

\section{Evidence, branch identification, and replay}\label{app:evidence}
The supplement preserves the numbered research packages as a fixed snapshot.
Their identifiers describe provenance, not additional unpublished theorems
which a reader must accept without a proof. The analytic arguments used
in this paper are given above; detailed local numerical witnesses and
coefficient differentiation formulas are in the indicated files.

\begin{center}\small
\begin{tabular}{p{.20\linewidth}p{.18\linewidth}p{.52\linewidth}}
\toprule
Claim & Evidence packages & Principal material\\
\midrule
General expansion & R21--R23 &
Finite and countable-switch proofs, exact algebra, examples\\
Cusp identity & R01, R03, R09, R29 &
Integral root box, connected uniqueness tubes, positive majorants,
32 driver cells\\
Theta coefficients & R24, R28, R29 &
Full root sums, weighted tails, implicit differentiation\\
Effective remainder & R25--R27, R30 &
Local jets, finite repair, moving prefix, scalar inversion\\
Actual-value order & R31 &
Moment identity, 1024 spatial cells, whole-tail comparison\\
\bottomrule
\end{tabular}
\end{center}

For reproducibility, the supplement includes the source, data, and dependency
closure of these packages. Its index maps each identifier to a directory,
manifest, proof, producer, and checker.
The following commands use paths relative to the supplement root:
\begin{verbatim}
python -B research/finite_slope_optimum/audit_snapshot.py
python -B research/finite_slope_fourth_order/audit_snapshot.py
python -B research/infinite_slope_fourth_order/audit_snapshot.py --with-arb
python -B research/cusp_order_transport/audit_snapshot.py --with-arb
\end{verbatim}
The last command recursively checks the later theta dependencies.
It regenerates the 1024 spatial interval enclosures, reconstructs the
32 driver-cell chain, and checks arithmetic and recorded prerequisites.
The accompanying report distinguishes freshly evaluated transcendental
enclosures from frozen prerequisite inputs.
The implementation uses Python, python-flint/Arb \cite{arb,integration},
exact rational arithmetic, and mpmath for non-rigorous diagnostics.
Pinned versions and input hashes are supplied.

Some details concerning identity and differentiation deserve emphasis.
The initial cusp is enclosed by an integral-based contraction, with
\(D_3D_4\ne0\). In the original continuation, the scaled contraction
bound is below \(0.502\) and the forcing-plus-contraction bound
is below \(0.560\).
At adjacent faces a smaller root enclosure lies inside the next full
uniqueness box. On the short interval used here, refined anchors and
degree-five parameter expansions with sixth-order remainder remain
inside those already identified tubes.
Thus a collection of unrelated zeros is not substituted for a branch.

The lower sign-moment equations have derivative \(-G\).
Their implicit solution \(b^*(\nu)\), and the simple residual zeros,
are differentiated together. For example,
\[
 \delta_0'=(f'D-fD')/D^2,\qquad
 C_2'=C_2(3f'/f-4D'/D+\Gamma'/\Gamma).
\]
The \(b^{*\,\prime}\) contribution to \(D'\) vanishes by the sign
moments; it does not vanish from \(\Gamma'\).
Root-sum derivatives retain root motion. Their tail bounds are justified
by separated zeros, implicit derivative bounds, and double exponential
majorants of the same type as \eqref{eq:weighted}.
The package R29 gives the termwise formulas and enclosing intervals
used for \eqref{eq:coefficientmotion}.

The final transport audit records 31 exact algebra groups,
19,017 rational decisions, and 1024 closed spatial cells. Its independent
high-precision integral and composition diagnostics are useful consistency
checks, but are not interval proofs.
The rational checker uses stated transcendental enclosures as inputs;
it is not a second independent rigorous integration engine.
The four replay entry points passed on 23 September 2026, and
1,334 entries in 32 pre-existing manifests remained unchanged.
These counts establish the tested snapshot and the coverage of that
replay. They do not establish that every line of analytic reasoning
is error-free.

For the public supplement, local machine paths in bookkeeping records are
replaced by neutral placeholders. All dependent checksum references are
updated consistently, and the complete transformation map is supplied.
No mathematical formula or numerical enclosure is changed by that
transformation. The public copy is replayed separately; its manifest
identities are therefore distinguished from the original local snapshot.
